\chapter{Introduction}\label{ch:introduction}

\section{General Description}\label{sec:general}
This report documents the modelling, identification and position control of a
DC motor that drives a pointer via a gearbox with belt drive (Box~1 of the
laboratory test benches). The controlled variable is the angular position of the
pointer, measured on the output shaft by an optical quadrature encoder. The motor
is supplied by a PWM-controlled motor driver, and all signals are processed by a
data acquisition card with MATLAB/Simulink and Simulink Desktop
Real-Time~\cite{laboranleitung}.

\section{Motivation}\label{sec:motivation}
DC motors are widely used electro-mechanical actuators, for example in
automotive applications, medical devices and industrial
automation~\cite{laboranleitung}. An exact positioning of the driven load
requires a closed-loop controller, whose design relies on a sufficiently
accurate model of the plant. Since not all parameters of a real drive train are
known, the model is derived both analytically and from measured data.

\section{Learning Goal}\label{sec:goal}
The aim of the laboratory is to identify an electro-mechanical system from its
measured step response and to design a suitable controller~\cite{laboranleitung}.
The step response is approximated by a second-order lag element, a PID
controller is designed with the Ziegler--Nichols method, and the closed loop is
simulated, implemented on the real test bench and evaluated.
